\documentclass[11pt]{article}
\usepackage[margin=1in]{geometry}
\usepackage[T1]{fontenc}
\usepackage{lmodern}
\usepackage{enumitem}

\setlist[itemize]{leftmargin=*, itemsep=0.25em, topsep=0.25em}
\setlist[enumerate]{leftmargin=*, itemsep=0.45em, topsep=0.35em}

\begin{document}

\section*{Executive Summary}
This paper remains a strong, publishable-quality contribution with a clear core advance: a shared empirical benchmark that applies competing operational definitions of life to identical traces under controlled ablations. The manuscript is broad in scope, technically careful, and mostly coherent from motivation to interpretation.

After a final paper-wide consistency check, the main risk is not missing experiments but claim calibration and cross-section alignment. The quantitative evidence is generally solid, yet a few high-level statements overreach relative to scope (especially substrate generalization and predictive ``best'' language), and key caveats are introduced later than their interpretive importance warrants.

The highest-impact revision path is editorial and structural: tighten evidence-weighted tone, stabilize terminology across abstract/results/discussion, make equation-level assumptions easier to audit, and foreground the strict-mode interpretation early. With these adjustments, the paper's narrative coherence and credibility would match its methodological ambition.

\section*{Prioritized Findings}
\textbf{The findings below are ranked by confidence (highest to lowest confidence).}

\begin{enumerate}
  \item \textbf{High confidence --- The benchmark contribution is clear and novel.} Evaluating D1--D4 on shared data with family ablations directly addresses a longstanding comparability gap in ALife.

  \item \textbf{High confidence --- The paper-wide narrative is mostly coherent.} The D2--D3 disagreement axis is introduced, tested, and interpreted consistently across major sections.

  \item \textbf{High confidence --- Predictive interpretation must be anchored on strict mode.} The D4 shift from legacy AUC 0.78 to strict AUC 0.54 is large and should be presented as a central qualifier, not a late robustness detail.

  \item \textbf{High confidence --- D1's structure creates a fairness-interpretation tension.} Geometric aggregation and cross-family $\beta$ dependence make D1 intentionally strict and context-sensitive, so ``best definition'' language should be scoped to the selected target and setup.

  \item \textbf{Medium-high confidence --- Statistical rhetoric needs tighter confirmatory/exploratory separation.} Some sentences read as definitive despite relying on descriptive intervals or uncorrected multi-comparison contexts.

  \item \textbf{Medium-high confidence --- Equation-level assumptions are under-specified in prose.} D1's $\sigma(\beta)$ transform, D2's Price-equation trait meaning, and D3's correction-order workflow need explicit definitions to improve auditability.

  \item \textbf{Medium confidence --- Terminology drifts across sections.} ``Unseen conditions'' is correctly defined as held-out seeds, but the phrase can still imply unseen regimes unless repeatedly scoped.

  \item \textbf{Medium confidence --- Lenia evidence is supportive but preliminary.} The observed D3 $>$ D2 pattern is useful cross-substrate signal, yet current breadth is best framed as early external support rather than substrate-invariant confirmation.

  \item \textbf{Medium confidence --- Reproducibility intent is strong but review-time status is unclear.} Artifact placeholders leave uncertainty about what a reviewer can execute immediately during blind review.
\end{enumerate}

\section*{Section Recommendations}

\subsection*{Abstract and Introduction}
\begin{itemize}
  \item Keep the current motivation and contribution statement, but soften strongest verbs so claims are evidence-weighted rather than definitive.
  \item Reiterate once in the abstract and once in results lead-in that ``unseen conditions'' means held-out seeds (100--199) within fixed regimes.
  \item Preview strict-mode implications in the introduction so later predictive findings do not read as a reversal.
\end{itemize}

\subsection*{Related Work}
\begin{itemize}
  \item Add one explicit contrast sentence separating conceptual pluralism from empirical comparability.
  \item Clarify that adapter-level comparability is methodological and does not collapse ontological differences among life concepts.
\end{itemize}

\subsection*{System and Baseline}
\begin{itemize}
  \item Keep architecture and mode choices, but move low-level implementation constants to supplementary material for readability.
  \item Use one stable unit-of-analysis phrase (regime $\times$ seed $\times$ family) across methods and results to reduce terminology drift.
  \item State briefly whether any simulator settings were adjusted after calibration inspection.
\end{itemize}

\subsection*{Definition Adapters (D1--D4)}
\begin{itemize}
  \item Define every transformation used in scoring equations (especially $\sigma(\beta)$) and give ranges so readers can reconstruct scores.
  \item For D2, clarify that $\Delta\bar{z}=\mathrm{Cov}(w,z)/\bar{w}$ with parent--child genome distance as $z$ measures selection on change magnitude, not direct ecological adaptation.
  \item For D3, present the inferential sequence as ordered steps (lag handling, TE/Granger decision, per-pair correction, FDR) so the final error-control claim is auditable.
  \item For D4, tie the causal weighting choice directly to strict-mode sensitivity and report whether moderate weight changes preserve ranking.
\end{itemize}

\subsection*{Experiments and Statistics}
\begin{itemize}
  \item Consolidate threshold usage in one table note: default vs calibrated, and which version is used in each figure/result block.
  \item Label confirmatory and exploratory analyses explicitly in subsection headings or captions for faster interpretive parsing.
  \item Keep AUC pairwise intervals, but add one short caution that family-wise error is not formally controlled across six comparisons.
\end{itemize}

\subsection*{Results}
\begin{itemize}
  \item Lead predictive-validity interpretation with strict-mode outcomes, then present legacy results as complementary context.
  \item Pair each headline quantitative claim with a practical-effect interpretation sentence, not only interval reporting.
  \item Maintain the D2--D3 disagreement axis as the qualitative center, while scoping all ranking language to target and substrate.
\end{itemize}

\subsection*{Discussion and Lenia}
\begin{itemize}
  \item Keep the conceptual synthesis framing, but describe Lenia findings as hypothesis-consistent external support rather than definitive invariance.
  \item Distinguish conclusions that are robust to operational variants from those that depend on scoring design choices.
  \item End with two concrete next tests: broader non-reproductive substrates and longer-run closure stability analysis.
\end{itemize}

\subsection*{Artifact and Presentation}
\begin{itemize}
  \item Provide a precise blinded reproducibility statement listing exactly what reviewers can run now and expected outputs.
  \item Ensure every quantitative figure caption includes sample size, split identity, and threshold policy (default/calibrated).
  \item Add a compact appendix roadmap linking each robustness claim to its validating figure or table.
\end{itemize}

\subsection*{Global Consistency}
\begin{itemize}
  \item Use one paper-wide claim template: ``within this substrate, regimes, and targets'' to avoid scope drift between sections.
  \item Keep terminology stable for key contrasts (legacy vs strict, confirmatory vs exploratory, disagreement vs prediction).
  \item Check that each strong conclusion appears only where the corresponding caveat is adjacent in the same subsection.
\end{itemize}

\end{document}